\subsection{Entornos Personalizados}

% --- Intro ---
\begin{frame}{Entornos Personalizados}
    El template incluye varios entornos predefinidos para estructurar las
    diapositivas de forma clara y visual.
\end{frame}

% --- Definiciones y Teoremas ---
\begin{frame}{Definiciones y Teoremas}
    \begin{DefinitionP}[title=Definición Formal]
        Utiliza este entorno para introducir formalmente nuevos conceptos,
        estructuras de datos o términos clave. El argumento opcional
        \texttt{title} permite dar un nombre descriptivo a la definición.
    \end{DefinitionP}

    \begin{TheoremP}[title=Propiedad Matemática]
        El entorno \texttt{TheoremP} sirve para destacar propiedades matemáticas
        fundamentales, lemas o teoremas que justifican la corrección o
        complejidad de un algoritmo.
    \end{TheoremP}
\end{frame}

% --- Consejos y Advertencias ---
\begin{frame}{Consejos y Advertencias}
    \begin{TipP}
        Este entorno se usa para dar consejos o recomendaciones prácticas que
        no son necesariamente reglas estrictas, pero ayudan a resolver
        problemas de manera más eficiente.
    \end{TipP}

    \begin{AlertP}
        Úsalo para advertir sobre errores comunes, desbordamientos de enteros
        (overflow), casos borde complicados o complejidad temporal excesiva.
    \end{AlertP}
\end{frame}

% --- Ejemplos ---
%% NOTA: [fragile] es necesario para incluir verbatim
\begin{frame}[fragile]{Ejemplos}
    \begin{ExampleP}[title=Caso de Prueba]
        Este cuadro es ideal para mostrar ejemplos concretos, como casos de
        entrada y salida, trazas de ejecución paso a paso, o diagramas
        explicativos.
        
        \textbf{Input:}
        \begin{verbatim}
        5
        1 2 3 4 5
        \end{verbatim}

        \textbf{Output:}
        \begin{verbatim}
        15
        \end{verbatim}
    \end{ExampleP}
\end{frame}
